\title{Parameter-Efficient Orthogonal Finetuning via Factorization}

\begin{abstract}
The prevalence of large foundation models has grown rapidly, but building these models from the ground up remains financially infeasible. This makes the development of strategies for adapting pre-existing models to new tasks, in a resource-efficient manner, a pressing concern. In this work, we focus on Orthogonal Finetuning (OFT), a principled approach for adapting foundation models to downstream tasks. Although OFT offers robust generalizability, it suffers from parameter inefficiency due to the high-dimensional nature of orthogonal matrices. To mitigate this, we first analyze OFT through the lens of information transmission and outline several core properties that promote superior parameter efficiency. Drawing inspiration from the Cooley-Tukey fast Fourier transform algorithm, known for facilitating efficient information transfer, we introduce a novel orthogonal parameterization based on butterfly structures. Integrating this with OFT yields Orthogonal Butterfly~(BOFT), a new approach to parameter-efficient finetuning. Notably, BOFT generalizes OFT by subsuming it as a specific case, offering a unified orthogonal finetuning framework. We further validate the efficacy of our method through comprehensive experiments involving the adaptation of large vision transformers, substantial language models, and text-to-image diffusion models for a wide range of tasks in vision and language domains.
\end{abstract}

\section{Introduction}

Modern architectures such as ChatGPT~\cite{brown2020language,chen2021evaluating} and Stable Diffusion~\cite{rombach2022high} highlight the outstanding ability of large foundation models to generalize across domains. This dramatic improvement in performance has been accompanied by exponential growth in parameter count (\emph{e.g.}, GPT-3 includes approximately 175B parameters). Consequently, training these foundational models entirely from scratch is becoming ever more challenging, both technically and economically. The field’s advancement now hinges on effectively repurposing these models for novel tasks without full retraining. To accomplish this, methods for task adaptation must be highly efficient. There are three predominant approaches: \emph{model finetuning}~\cite{hulora2022,zaken2022bitfit,guo2020parameter,raffel2020exploring,qiu2023controlling,zhang2022adaptive,chavan2023one}, where select parameters are optimized while maintaining the model’s original structure; \emph{adapter tuning}~\cite{rebuffi2017learning,houlsby2019parameter,pfeiffer2020adapterfusion,lin2020exploring,he2021towards}, where new trainable modules are introduced; and \emph{prompt tuning}~\cite{li2021prefix,lester2021power}, which attaches tunable prefix tokens to model inputs. Among these, model finetuning stands out due to its conceptual simplicity, impressive effectiveness, and lack of added inference delay.

The core concept of model finetuning is to maintain a close resemblance between the pretrained and finetuned models according to specific metrics, thereby safeguarding the knowledge acquired during pretraining. Typically, prevailing finetuning techniques utilize a small learning rate, an ad hoc tactic that limits the deviation between the initial and adapted model. Given that weight matrices have inherent structural properties, more methodical strategies aim to maintain the relationships within these matrices—for example, by preserving the pairwise angular relationships between neurons. Building on this perspective, the Orthogonal Finetuning~(OFT)~\cite{qiu2023controlling} framework has been proposed. In OFT, every neuron within a layer undergoes transformation via a shared orthogonal matrix, ensuring the theoretical conservation of neuron-to-neuron angles throughout the finetuning stage. While OFT has yielded strong results in terms of generalizability and convergence when applied to text-to-image diffusion models~\cite{qiu2023controlling}, its reliance on high-dimensional orthogonal matrices can make it parameter-intensive. To mitigate this, OFT implements a block-diagonal structure, effectively lowering the number of trainable parameters. Nonetheless, this efficiency incurs certain trade-offs—the block-diagonal structure enforces a fixed sparsity pattern, applying orthogonal transformations independently across distinct blocks. Although this arrangement saves parameters, it introduces potentially problematic inductive biases (\emph{e.g.}, it restricts expressiveness and precludes the ability to represent general linear transformations). Moreover, optimal selection of a sparsity pattern remains an open problem.

Addressing this challenge hinges on creating a dense orthogonal matrix while maintaining parameter efficiency. Though a dense orthogonal matrix of dimension $d$ conventionally demands $\mathcal{O}(d^2)$ parameters, our approach innovates by assembling such a matrix from a product of multiple factorized sparse matrices. This method relies on the observation that fewer trainable parameters can be used by opting for increased computational overhead rather than expanded storage. Since our construction relies on repeated multiplication of sparse matrices, these multiplications are performed at each training step. Conceptually, this factorization can be reframed as an information transmission scenario: forming a dense orthogonal matrix via the product of sparse matrices parallels information dissemination across a grid-based graph. This viewpoint allows for the design of various sparse matrix factorizations that constrain parameter counts while retaining sufficient expressivity to produce dense matrices.

Within the information transmission perspective, we achieve parameter efficiency by drawing upon the butterfly structures characteristic of the Cooley-Tukey fast Fourier transform algorithm~\cite{cooley1965algorithm}, known for their efficient information exchange across nodes~\cite{klasing1994broadcasting}. The butterfly graph naturally yields a sparse matrix decomposition referred to as butterfly factorization. Utilizing butterfly factorization enables the construction of a $d$-dimensional dense matrix as the product of $\mathcal{O}(\log d)$ sparse matrices, each containing only $\mathcal{O}(d)$ non-zero elements. By ensuring each sparse factor maintains orthogonality, this technique permits an orthogonal parameterization requiring just $\mathcal{O}(d\log d)$ parameters—substantially fewer than traditional approaches. Leveraging this efficient orthogonal structure, we introduce a new parameter-efficient finetuning technique—Orthogonal Butterfly~(BOFT). BOFT generalizes the OFT method as a specific instance within its broader orthogonal finetuning schema. Both the block-diagonal and butterfly constructions share an emphasis on sparsity, imposing a structural constraint that limits the number of trainable parameters in orthogonal matrices. Comparing our method to the low-rank approach in LoRA~\cite{hulora2022}, it is worth noting that both low-rank and sparse matrices belong to foundational families of structured matrices~\cite{candes2011robust} that aim for parameter efficiency.

In contrast to the block-diagonal structure used in OFT, which negotiates between expressiveness and regularity, BOFT leverages the butterfly structure to achieve a more gradual interpolation between full orthogonal matrices (enabling expressivity) and the identity matrix (encouraging regularity). This approach facilitates the identification of a more compact hypothesis space within the orthogonal group for various downstream applications. Given the butterfly structure’s extensive deployment in rapid linear operations such as the discrete Fourier and cosine transforms, we posit that our structurally-motivated strategy for parameter efficiency embeds an implicit inductive bias, which can enhance generalizability and help mitigate overfitting. This perspective arises from the butterfly structure's capacity to capture a broad class of well-known linear transformations, a feat that the block-diagonal structure in OFT is fundamentally incapable of matching. Our main contributions are summarized as follows:

\begin{itemize}[leftmargin=*,nosep]
    \item We address parameter efficiency in orthogonal finetuning via a novel information transmission perspective, reformulating the challenge of designing a parameter-efficient dense orthogonal matrix as an information flow problem on a grid-structured graph.
    \item Drawing inspiration from the butterfly structures underpinning the Cooley-Tukey algorithm, we introduce Orthogonal Butterfly, an orthogonal finetuning technique that is highly parameter-efficient and operates within our information transmission paradigm.
    \item We offer theoretical analyses explaining how BOFT dramatically reduces the number of trainable parameters, yet retains substantial expressive power and training stability. Owing to its matrix factorization, BOFT also enables an appealing form of weight interpolation (see Figure~\ref{fig:boft_interpolation}).
    \item We are the first to systematically deploy orthogonal finetuning across a spectrum of tasks beyond controllable text-to-image generation~\cite{qiu2023controlling}, thereby establishing its promise as a general-purpose finetuning mechanism. Specifically, we evaluate BOFT on a variety of adaptation problems spanning computer vision and natural language processing. BOFT demonstrates clear improvements over contemporary leading techniques, substantiating its strong parameter efficiency and broad generalization capabilities.
\end{itemize}

\textbf{Parameter-efficient finetuning (PEFT).} The escalation in scale and capability of foundation models (\emph{e.g.}, \cite{brown2020language,radford2021learning,rombach2022high,kirillov2023segment}) has fueled significant research into efficient adaptation of these models for specific tasks, with a focus on minimizing the number of parameters trained~\cite{treviso2023efficient,ding2023parameter,lialin2023scaling}. Among the diverse approaches proposed within the PEFT literature~\cite{houlsby2019parameter,aghajanyan2020intrinsic,hulora2022,edalati2022krona,wang2022adamix,gheini2021cross,zaken2022bitfit,guo2020parameter,sung2021training,ansell2022composable,lester2021power,li2021prefix,vu2022spot,he2021towards,mao2021unipelt,karimi2021compacter,liu2022few,sung2022lst,chen2022parameter,jia2022visual,chen2022adaptformer,zhang2022neural,jie2023fact,lian2022scaling,luo2023towards}, those relying on reparameterization strategies~\cite{aghajanyan2020intrinsic,hulora2022,edalati2022krona} are most closely related to our contributions. A prominent technique, LoRA~\cite{hulora2022}, modifies pretrained weights by incorporating the product of two low-rank matrices and demonstrates strong outcomes on a variety of language understanding tasks. As LoRA enforces a fixed rank for all added matrices, subsequent approaches~\cite{zhang2022adaptive,valipour2022dylora,zhang2023increlora} introduce mechanisms to dynamically modulate the rank per layer, thereby optimizing the allocation of parameters across the network. Owing to its effectiveness and ease of implementation, this low-rank adaptation strategy enjoys widespread adoption~\cite{dettmers2023qlora,zi2023delta,chavan2023one}. Building upon insights from the role of hyperspherical energy in supporting generalization~\cite{liu2018learning,liu2021orthogonal}, \cite{qiu2023controlling} put forth orthogonal finetuning (OFT) as a novel and potent alternative for text-to-image diffusion model adaptation. In OFT, a learnable orthogonal transformation is applied within each layer, leading to improved generalization capabilities and typically more stable optimization compared to LoRA. Nevertheless, OFT tends to require a greater number of trainable parameters. Enhancing the parameter-efficiency of OFT thus presents an appealing direction. Furthermore, the extent to which OFT generalizes to broader classes of adaptation problems—beyond its initial success in text-to-image diffusion~\cite{qiu2023controlling}—remains unexplored. To address this, BOFT leverages butterfly factorization to improve OFT's parameter utilization. This advancement enables the demonstration of orthogonal finetuning's utility across a wider array of adaptation scenarios.

\textbf{Butterfly structures.} The radix-2 Cooley-Tukey procedure systematically decomposes the $N$-point discrete Fourier transform into two smaller transforms of size $\frac{N}{2}$, giving rise to a distinctive butterfly pattern which is mathematically represented by the product of several sparse matrices (often referred to as a butterfly matrix). Such matrices have been employed for orthogonal matrix parameterizations to circumvent the need for pivoting in Gaussian elimination and enhance computational efficiency~\cite{parker1995random}; they have also proved effective in improving the robustness of recurrent neural network training~\cite{jing2017tunable} and in the context of kernel approximation~\cite{munkhoeva2018quadrature}. Recent work~\cite{dao2019learning,dao2019kaleidoscope} demonstrates learning rapid linear transforms through butterfly matrix representations, while \cite{chen2022pixelated,dao2022monarch} exploit these structures to train neural networks more efficiently. Beyond these applications, butterfly configurations facilitate fast matrix-vector multiplications~\cite{michielssen1996multilevel,o2010algorithm}, enable efficient approximation of data-sparse matrices~\cite{li2015butterfly}, and support various network communication protocols~\cite{klasing1994broadcasting,li2003linear,rajasingh2016transmission}. Distinct from this prior literature, our approach targets the application of butterfly structures to substantially improve the parameter-efficiency of OFT.

\section{An Information Transmission View on Orthogonal Finetuning}

We begin by outlining foundational concepts of OFT. To adapt a pretrained weight matrix, OFT reformulates the updated weight matrix as the product of a trainable orthogonal matrix and the fixed pretrained weight matrix. In contrast with LoRA—which modifies weights by adding a learnable low-rank matrix—OFT applies a multiplicative orthogonal matrix for the update. LoRA achieves efficiency by employing low-rank parameterization; in comparison, traditional OFT utilizes a block-diagonal orthogonal matrix~\cite{qiu2023controlling}, where decreasing block sizes lead to higher parameter efficiency. Figure~\ref{fig:comp_lora} provides an intuitive contrast between the two. The rationale behind leveraging orthogonal transformations for finetuning is to retain pair-wise angles among neurons~\cite{liu2018learning,liu2021orthogonal,qiu2023controlling}, thereby conserving much of the pretrained model’s semantic information.

In practical terms, OFT learns an orthogonal matrix $\bm{R}\in\mathbb{R}^{d\times d}$ for a given pretrained linear layer $\bm{W}^0\in\mathbb{R}^{d\times n}$, updating the forward operation from $\bm{z}=(\bm{W}^0)^\top\bm{x}$ to $\bm{z}=(\bm{R}\bm{W}^0)^\top\bm{x}$, with $\bm{x}\in\mathbb{R}^{d}$ as input and $\bm{z}\in\mathbb{R}^{n}$ as output. OFT uses the identity matrix to initialize $\bm{R}$, which realizes a zero initialization. To guarantee that $\bm{R}$ remains orthogonal during finetuning, we follow \cite{qiu2023controlling,liu2021orthogonal} by using the Cayley transform for parameterization: specifically, $\bm{R}=(\bm{I}+\bm{Q})(\bm{I}-\bm{Q})^{-1}$ where $\bm{Q}$ is skew-symmetric, i.e., $\bm{Q}=-\bm{Q}^\top$. For greater efficiency, the block-diagonal scheme structures $\bm{R}$ as $\text{diag}(\bm{R}_1,\bm{R}_2,\cdots,\bm{R}_r)$, with each block $\bm{R}_i\in\mathcal{R}^{b\times b}$ being an orthogonal matrix and $br=d$. While this block-diagonal design yields parameter savings, it implicitly assumes that the neuron's dimension (i.e., columns of $\bm{W}^0$) is partitioned into $r$ segments, with each group processed independently via separate orthogonal matrices. Despite the practical successes of the block-diagonal approach, arbitrarily dividing neuron dimensions by index offers no principled justification, pointing to the appeal of dense orthogonal matrices. This naturally leads to the central question: \emph{Is it feasible to design a dense orthogonal matrix that preserves parameter efficiency?}

To tackle this issue, we introduce a method that constructs a dense orthogonal matrix by multiplying several sparse orthogonal matrices together. To offer a comprehensive yet accessible framework for analyzing the sparsity structures involved in orthogonal matrix factorization, we reinterpret the task of creating a dense orthogonal matrix within OFT as analogous to an information transmission process. In particular, the generation of a dense matrix $\bm{R}\in\mathbb{R}^{d\times d}$ as a product of $m$ square matrices, denoted $\thickmuskip=2mu \medmuskip=2mu \bm{R}=\bm{B}_m\bm{B}_{m-1}\cdots\bm{B}_1$, can be modeled as transmitting signals across a lattice composed of $d\times (m+1)$ nodes (see Figure~\ref{fig:info_trans}). The rationale for adopting the information transmission lens stems from viewing a $d$-by-$d$ dense matrix as representing full connections between two groups of $d$ nodes each. Within the matrix $\bm{R}$, if $R_{ij}$ equals zero, this implies that no information flows from the $j$-th to the $i$-th node, whereas a nonzero entry signifies a possible transmission path. Thus, expressing $\bm{R}$ as a product $\bm{B}_m\bm{B}_{m-1}\cdots\bm{B}_1$ equates to executing sequential message passing as defined by the graph structures induced by each $\bm{B}_i$. The process initiates transmission through $\bm{B}_1$, progressing in order to $\bm{B}_n$. 

Figure~\ref{fig:info_trans} provides a concrete scenario using the factorization $\bm{R}=\bm{B}_5\bm{B}_4\bm{B}_3\bm{B}_2\bm{B}_1$, with associated sparsity structures and induced graphs illustrated. Here, the depicted graph arises from systematically expanding the matrix multiplication over time. Each matrix $\bm{B}_i$ serves as an adjacency matrix encoding links between the $i$-th and $(i+1)$-th levels; specifically, the $(j_1,j_2)$ entry of $\bm{B}_i$ indicates whether a directed connection exists from the $j_2$-th node at level $i$ to the $j_1$-th node at level $i+1$ (where a zero indicates no connection). To achieve a dense matrix through $\bm{B}_5\bm{B}_4\bm{B}_3\bm{B}_2\bm{B}_1$, it is necessary that each node at the starting level can communicate with all nodes at the $6$-th level. Considering only the factorization $\bm{R}=\bm{B}_4\bm{B}_3\bm{B}_2\bm{B}_1$, which involves source nodes in the initial layer and receiver nodes at the $5$-th layer, one observes that, for instance, node $1$ fails to transmit information to node $3$, meaning the corresponding position $(3,1)$ in the sparsity pattern is zero. This correspondence between the graphical structure and the zero-pattern similarly manifests in factorizations like $\bm{R}=\bm{B}_3\bm{B}_2\bm{B}_1$ and $\bm{R}=\bm{B}_2\bm{B}_1$. More broadly, for a $d\times d$ matrix $\bm{R}$ to exhibit full density, the collection of $m$ factor matrices $\bm{B}_m,\ldots,\bm{B}_1$ must collectively specify a system of directed edges over a $d\times (m+1)$ network, where each edge links two nodes in adjacent layers (that is, columns), ensuring that all source nodes in the first layer are able to convey information to every destination node in the $(m+1)$-th layer.

The information transmission framework offers valuable insights into the sparsity structures that emerge in factorization matrices within OFT. As depicted in Figure~\ref{fig:blk_diag}, the original OFT employs a block-diagonal formation in $\bm{R}$. While this configuration effectively decreases the number of trainable parameters, it does not permit the realization of a fully dense matrix $\bm{R}$. The primary objective is thus to assemble a dense orthogonal matrix by integrating $m$ sparse orthogonal matrices, aiming to minimize the number of necessary trainable parameters. Within the context of the information transmission perspective, two main criteria are outlined: \emph{(i)} \textbf{dense connectivity}, whereby every node in the input layer can reach every node in the output layer through at least one path, and \emph{(ii)} \textbf{minimum free edges}, stipulating that, given the orthogonality constraint, the total edge count should be minimized. Orthogonality imposes nuanced constraints on the connections between adjacent layers. Specifically, to maintain full-rank in each matrix $\bm{B}_i$—an essential condition for orthogonality—it is necessary to have $d$ edges connecting all nodes from the $i$-th layer to the $(i+1)$-th layer, ensuring a bijective mapping. This necessitates a minimum of $d$ inter-level edges (\emph{e.g.}, 4 as shown in Figure~\ref{fig:info_trans}); however, these edges are structurally determined by orthogonality and do not represent trainable parameters (for example, in a $d\times d$ orthogonal matrix with $d$ nonzero entries, each entry must be $\pm 1$). As a result, orthogonal matrices inherently use fewer parameters than unrestricted matrices, introducing extra dependencies among the edges. For instance, in a $2\times 2$ orthogonal matrix, a zero at position $(1,1)$ necessitates a zero at $(2,2)$ (\emph{i.e.}, removing one edge may require the removal of another). Consequently, the orthogonality condition can alter the set of feasible edge configurations by implicitly adding or removing edges. By examining the pattern of nonzero elements in the composite orthogonal matrix, the information transmission viewpoint is especially pertinent to OFT, as our main concern lies with the distribution of trainable nonzero elements in $\bm{R}$, rather than their precise values.

A straightforward fully-connected mapping between two hierarchical levels entails $\mathcal{O}(d^2)$ connections (\emph{i.e.}, a single dense orthogonal matrix), specifically resulting in $d^2-d$ distinct edges (for instance, when $d=4$, there are 12 such connections). As shown in Figure~\ref{fig:info_trans}, one can construct a matrix factorization with only $10$ edges in total, which is even fewer than the connections required for a complete dense orthogonal transformation. This approach allows an investigation of OFT's parameter efficiency through the lens of graph structures, providing an intuitive platform for devising alternative factorizations. Drawing upon structural insights from the Cooley-Tukey algorithm, we note that the so-called butterfly graphs~\cite{cooley1965algorithm} efficiently interlink $d$ input nodes with $d$ output nodes using merely $\mathcal{O}(d\log d)$ connections. To illustrate, the configuration in Figure~\ref{fig:info_trans} realizes full connectivity with $10$ edges, whereas the butterfly topology accomplishes the same with just $8$. In the following, we detail how adopting the butterfly architecture contributes to improved parameter efficiency.

\section{Orthogonal Parameterization by Butterfly Factorization}

The butterfly arrangement was originally introduced within the Cooley-Tukey algorithm to accelerate the Fourier transform computation. In this context, a slight modification in the frequency domain induces a widespread impact across the spatial domain, mirroring our own scenario in information transmission: every node at the initial stage can convey data to every node in the terminal stage. Beyond signal processing, butterfly configurations have been employed extensively as an effective architecture for computer networking~\cite{solihin2015fundamentals,li2011linear} to support rapid and scalable information flow. Let us consider $k \geq 2$ as an integer power of $2$; we define the \emph{butterfly factor} $\bm{B}^F(k)$ by

\begin{equation}\label{eq:but_factor}
\begin{aligned}
    \bm{B}^F(k)=\begin{bmatrix}
    \text{diag}(\bm{d}_1) & \text{diag}(\bm{d}_2) \\
    \text{diag}(\bm{d}_3) & \text{diag}(\bm{d}_4)
    \end{bmatrix}\in\mathbb{R}^{k\times k},
\end{aligned}
\end{equation}

where each $\bm{d}_i \in \mathbb{R}^{\frac{k}{2}},\ \forall i$, denotes an arbitrary vector. With $d=2^N$, the $d$-dimensional \emph{butterfly matrix} $\bm{B}(d) \in \mathbb{R}^{d\times d}$ is specified recursively as follows:

\begin{equation}
\begin{aligned}
    \bm{B}(d) = \tilde{\bm{B}}(d,d)\cdot\begin{bmatrix}
    \bm{B}_1(\frac{d}{2}) & \bm{0} \\
    \bm{0} & \bm{B}_2(\frac{d}{2})
    \end{bmatrix} = \tilde{\bm{B}}(d,d)\tilde{\bm{B}}(d,\frac{d}{2})\cdots\tilde{\bm{B}}(d,2),
\end{aligned}
\end{equation}

where $\bm{B}_1(\frac{d}{2})$ and $\bm{B}_2(\frac{d}{2})$ are butterfly matrices of dimension $\frac{d}{2}$. The \emph{butterfly component} is subsequently introduced as $\thickmuskip=2mu \medmuskip=2mu \tilde{\bm{B}}(d,k)=\text{diag}(\bm{B}^F_1(k),\cdots,\bm{B}^F_{d/k}(k))$, representing a block-diagonal matrix with overall shape $d\times d$ and block size $k$, where each $\bm{B}^F_i(k)$, for every $i$, is a butterfly factor described in Equation~\ref{eq:but_factor}. With this setup, we are prepared to leverage the butterfly matrix framework for orthogonal matrix parameterization. To facilitate this, it suffices that all block-diagonal factors $\tilde{\bm{B}}(d,k)$ comprising the butterfly matrix $\bm{B}(d)$ are themselves orthogonal. Consider first the specific case of $\tilde{\bm{B}}(d,2)$, which features $2 \times 2$ diagonal blocks; orthogonality of $\tilde{\bm{B}}(d,2)$ is readily attained by using the Cayley transform or a simple $2$-dimensional rotation for each block, following the methodology of \cite{liu2021orthogonal,qiu2023controlling}. Since the nonzero structure of any butterfly component corresponds, up to a permutation, to that of $\tilde{\bm{B}}(d,2)$, every butterfly component can therefore be parameterized to be orthogonal through analogous constructions. In this manner, one achieves an efficient way to parameterize orthogonal matrices using multiple $2\times2$ orthogonal matrices as building blocks. Extending the idea per \cite{chen2022pixelated}, we define a block butterfly component $\tilde{\bm{B}}^b(d,k)$ in which each element of $\bm{d}_i$ (for each $i$) is a $b\times b$ matrix rather than a scalar. To ensure the block butterfly component $\tilde{\bm{B}}^b(d,2)$ remains orthogonal, it is sufficient to parameterize each $2b\times 2b$ block as an orthogonal matrix. Since the nonzero structure of any $\tilde{\bm{B}}^b(d,k)$ for $k > 2$ can be interpreted as a block-wise permutation of $\tilde{\bm{B}}^b(d,2)$, these, too, can be consistently parameterized to ensure orthogonality. Bringing these elements together, the forward computation in BOFT becomes

\begin{equation}
    \begin{aligned}\nonumber
    \bm{z}=\big(\bm{R}(m,b)\cdot\bm{W}^0\big)^\top\bm{x},~~~~\text{s.t.}~\bigg{\{}\bm{R}(m,b)=\prod_{i=1}^m \tilde{\bm{B}}^b_{(d,i)}~~\&~~\underbrace{\big(\tilde{\bm{B}}^b_{(d,j)}\big)^\top\tilde{\bm{B}}^b_{(d,j)}=\tilde{\bm{B}}^b_{(d,j)}\big(\tilde{\bm{B}}^b_{(d,j)}\big)^\top\!=\bm{I}_d}_{\forall j\in [1, m]}\bigg{\}},
    \end{aligned}
\end{equation}

In this formulation, we use $\tilde{\bm{B}}^b(d,2^{m - i+1})$ as shorthand for $\tilde{\bm{B}}^b_{(d,i)}$, and $\bm{I}_d$ represents the $d$-dimensional identity matrix. The orthogonal matrix $\bm{R}(m,b)\in\mathbb{R}^{d\times d}$ arises as a product of multiple orthogonal butterfly matrices. We further denote BOFT instantiated with $\bm{R}(m,\frac{b}{2})$ as BOFT($m$,$b$), under the requirement that $b\geq 2$. For the specific setting where $\thickmuskip=2mu \medmuskip=2mu m=1$, BOFT($1$,$b$) specializes to the block-diagonal OFT~\cite{qiu2023controlling} featuring block size $b$. Setting $b=d$, BOFT($1$,$d$) recovers the standard OFT~\cite{qiu2023controlling} involving a fully unconstrained orthogonal matrix. By employing $\bm{B}^{\frac{b}{2}}(d)$ as the block butterfly matrix, BOFT($\log \frac{2d}{b}$,$b$) constructs $\bm{R}$ as a dense orthogonal matrix. In the general case, the parameter count for BOFT($m$,$b$) scales as $\frac{1}{2}(b-1)dm$, which represents the effective number of trainable parameters needed to adapt a linear transformation of shape $d\times n$. If we select the pure butterfly case, that is, $m=\log d, b=2$, then BOFT incurs a parameter complexity of $\mathcal{O}(d\log d)$. In comparison, a traditional OFT built from a dense orthogonal matrix needs $\mathcal{O}(d^2)$ parameters, and block-diagonal OFT with $r$ blocks requires $\mathcal{O}(bd)$. Thus, whereas the original OFT must utilize $b=d$ to achieve a dense orthogonal transformation, BOFT can flexibly produce dense orthogonal matrices using any suitable block size $b$.

\textbf{Identity initialization for BOFT}. It is common in finetuning practices to initialize from the exact pretrained weights, preserving close proximity to the original model during adaptation. For instance, LoRA employs zero initialization for the introduced low-rank parameters. In the case of BOFT, we opt to initialize each butterfly component as the identity (i.e., initializing the corresponding skew-symmetric matrix as zero in the Cayley transform), thereby starting from a configuration identical to the pretrained model.

\textbf{Multiplicative Dropout for BOFT}. In the context of LoRA~\cite{hulora2022}, Dropout regularization is employed to mitigate overfitting within the low-rank adaptation framework, leveraging standard Dropout~\cite{srivastava2014dropout}. However, this form of Dropout is not directly compatible with BOFT because of the multiplicative nature of the parameter updates. To bridge this gap, we introduce a new variant—multiplicative Dropout—designed specifically for BOFT. This approach leverages the modular structure of the orthogonal matrix $\bm{R}(m,b)$, which is built as a product of $m$ orthogonal butterfly factors, each of which can be organized into $2b\times 2b$ block-diagonal matrices. In our method, a random selection of $p_1$ percent of the butterfly factors and $p_2$ percent of block-diagonal entries within each factor are identified and replaced with identity matrices

\textbf{Expressivity of BOFT}. While the butterfly structure combined with permutations has been demonstrated to exactly reproduce several well-known fast linear transforms~\cite{dao2019learning,dao2019kaleidoscope} (\emph{e.g.}, fast Fourier transform, Hadamard transform), it remains an open question how effectively the orthogonal butterfly matrix can serve as an approximation to arbitrary orthogonal matrices. To investigate this, we carried out a numerical experiment aiming to approximate a randomly generated dense orthogonal matrix~\cite{anderson1987generation} of size $512\times 512$. In Figure~\ref{fig:para_eff}, the results are reported as averages across 10 different random seeds. Here, the y-axis reflects the approximation error, and the x-axis indicates the number of effective trainable parameters. For each color, the curves correspond to BOFT models configured with the same block size; the leftmost data point of each color shows the approximation error for BOFT($1$,$b$), which corresponds to the conventional block-diagonal OFT with block size $b$. Overall, BOFT exhibits greater parameter efficiency than OFT. For instance, BOFT($9$,$2$) achieves higher expressivity compared to BOFT($1$,$16$) while employing significantly fewer parameters. More broadly, BOFT with smaller block sizes ($b$) and larger butterfly depths ($m$) tends to use parameters more efficiently; for example, BOFT($6$,$4$) achieves an error comparable to BOFT($2$,$16$) but with a lower parameter count. This demonstrates that the butterfly structure encompasses a richer subset of the orthogonal group compared to block-diagonal structures, thus rendering BOFT provably more expressive than OFT for the same block size.

\begin{theorem}[Expressivity of BOFT]\label{thm:exp_boft}
    BOFT provides greater expressivity than OFT with equal block size. To enable the butterfly matrix to approximate any orthogonal matrix of size $d$, one can consider a product of the form $\bm{B}_{d-1,1}(d)\bm{B}^\top_{d-1,2}(d)\cdots\bm{B}_{1,1}(d)\bm{B}^\top_{1,2}(d)$, where each $\bm{B}_{i,j}(d),\forall i,\forall j$ is itself a butterfly matrix.
\end{theorem}

According to Theorem~\ref{thm:exp_boft}, BOFT can be naturally extended—where the final orthogonal matrix is constructed as $\thickmuskip=2mu \medmuskip=2mu \bm{R}^G(m_1,b_1,m_2,b_2,l)=\bm{R}_{l,1}(m_1,b_1)\bm{R}_{l,2}^T(m_2,b_2)\cdots\bm{R}_{1,1}(m_1,b_1)\bm{R}_{1,2}^T(m_2,b_2)$, with $\bm{R}_{i,j}^T(m,b)$ representing the orthogonal matrices within the BOFT framework. Notably, when $m_1=m_2=\log d$, $b_1=b_2=2$, and $l=d-1$, this formulation can encompass the entire orthogonal group, a structure often referred to as the kaleidoscope hierarchy~\cite{dao2019kaleidoscope}. Nevertheless, it is important to note that increased expressivity does not inherently guarantee improved finetuning performance; indeed, full finetuning—despite being universally expressive—may fail to achieve satisfactory results. Balancing the interplay between expressivity and regularization is crucial for ensuring generalization in model finetuning. By incorporating structural priors, BOFT expands the feasible parameter space, thereby enabling a more favorable balance between expressiveness and regularity.

\textbf{Spectral properties}. Orthogonal finetuning typically achieves superior spectral characteristics compared to LoRA, as it exactly maintains the spectral norm of the original pretrained matrix $\bm{W}^0$. This can be seen by considering the singular value decomposition: $\bm{W}^0=\bm{U}\bm{\Sigma}\bm{V}^\top$, with $\bm{U}$ and $\bm{V}$ orthogonal and $\bm{\Sigma}$ a diagonal matrix containing the singular values. Both OFT and BOFT perform a left multiplication by an orthogonal matrix $\bm{R}$, producing updated weights of the form $\bm{R}\bm{U}\bm{\Sigma}\bm{V}^\top$, which leaves the maximum singular value, and thus the spectral norm of $\bm{W}^0$, unchanged. The conservation of the spectral norm has been demonstrated to substantially improve both the stability of training and model generalization~\cite{miyato2018spectral,yoshida2017spectral}. Further noteworthy mathematical properties can be found in Appendix~\ref{app:sn_property}.

\textbf{Orthogonal finetuning as learning bilinear similarity}. BOFT may be interpreted as optimizing a bilinear similarity measure given by $\bm{w}^0_i\bm{R}\bm{x}$, where $\bm{w}^0_i$ denotes the $i$-th column of $\bm{W}^0$, corresponding to an individual neuron. In this context, BOFT is concerned with learning the matrix $\bm{R}$ for bilinear similarity, under the constraint that $\bm{R}$ is orthogonal. This formulation is intrinsically related to concepts in distance metric learning~\cite{xing2002distance} as well as bilinear forms~\cite{roman2005advanced}.

\textbf{Inductive bias and generalization in BOFT}. In BOFT, $\bm{R}(m,b)$ typically belongs to a specific subset of the orthogonal group, introducing restrictions to the set of permissible hypotheses and consequently embedding an inductive bias. We contend that the structured inductive bias imposed by butterfly factorization enhances generalization, since this structure mirrors the patterns found in many foundational linear transformations~\cite{dao2019kaleidoscope}, such as the discrete Fourier transform, discrete sine/cosine transform, and Hadamard transform. In addition, the use of sparse matrix factorization in BOFT potentially introduces further forms of implicit inductive bias~\cite{gunasekar2017implicit,li2018algorithmic,liu2019neural}.

\textbf{Comparison to butterfly-based sparse training}. Several prior studies~\cite{dao2019learning,dao2019kaleidoscope,chen2022pixelated,dao2022monarch} have investigated sparse training regimes employing the butterfly parameterization. Most of these works center on directly reparameterizing weight matrices with the butterfly structure and training neural architectures from the ground up. Notably, \cite{dao2022monarch} explores a finetuning methodology wherein pretrained weights are initially projected onto a variation of butterfly matrices, after which the resultant components are optimized for downstream applications. In contrast, BOFT introduces a novel finetuning approach, utilizing transformation of the weights via layer-wise shared weight matrices.

\input{rebuttal-results}

\section{Concluding Remarks and Limitations}

This work introduces Orthogonal Butterfly, a broadly applicable and parameter-efficient finetuning technique for foundation models, rooted in the butterfly architecture. The central concept for achieving higher parameter efficiency involves expressing a dense orthogonal matrix as the product of multiple sparse orthogonal matrices. To streamline the discovery of valid matrix factorizations, a graphical information transmission framework is devised. Within this perspective, it becomes evident that the butterfly structure is particularly suited for constructing sparse orthogonal matrix factorizations aligned with our objectives. Empirical results establish the effectiveness of BOFT for finetuning large-scale models, including large language models, substantial vision models, and text-to-image generation models. Furthermore, experiments underscore the advantages of BOFT as a general-purpose finetuning methodology.

However, despite promising empirical performance, BOFT does have some limitations. Since the resulting orthogonal matrix in BOFT is formed as a product of several orthogonal matrices, the training process incurs slightly higher runtime overhead compared to OFT. Enhancing the training efficiency of BOFT thus remains a subject for further exploration. Encouragingly, after finetuning is completed, the orthogonal matrices learned by BOFT can be fused directly into the original model, incurring no extra inference-time cost. Additionally, it is not yet clear whether the butterfly network represents the optimal architecture for information transfer. The proposed information transmission framework opens the possibility of drawing from insights in computer networking, where the efficacy of network topologies for communication is rigorously examined. This inspires the prospect that alternative, possibly more efficient network structures might be leveraged for constructing dense orthogonal matrices.